El presente trabajo abordó la reingeniería mecánica y la automatización de un brazo robótico cartesiano existente para su aplicación en la línea de empaque del Laboratorio de Diseño de Procesos de la Universidad del Valle de Guatemala. El sistema existente presentaba limitaciones mecánicas en los ejes X, Y y Z, entre ellas fricción elevada, desalineaciones y movimientos irregulares, las cuales afectaban su precisión, repetibilidad y capacidad de operar de forma estable. Para corregir estas deficiencias, se rediseñaron parcialmente los componentes estructurales y los sistemas de traslación mediante herramientas de diseño asistido por computadora. Asimismo, se incorporaron ejes guía y rodamientos lineales con el propósito de mejorar la fluidez y estabilidad del movimiento.

La automatización del sistema se desarrolló mediante una arquitectura de control embebido encargada de coordinar los motores, sensores, dispositivos de comunicación y periféricos. Se integraron finales de carrera para realizar la calibración automática y establecer referencias de movimiento, así como un codificador rotativo para medir la velocidad de la banda transportadora y facilitar la sincronización del proceso. Además, se implementó un sistema de visión computacional para identificar objetos con características previamente definidas, como figura y color, y determinar la información necesaria para ejecutar tareas de detección, clasificación, manipulación y transferencia mediante una secuencia de tipo \textit{pick and place}. El brazo también contó con un modo de operación manual mediante un mando inalámbrico y una interfaz local para visualizar el estado de calibración, el modo de operación y las acciones realizadas por el sistema.

La metodología se organizó en fases de evaluación del mecanismo existente, rediseño mecánico, fabricación y ensamblaje, desarrollo de la arquitectura electrónica, programación del control de movimiento, integración del sistema de visión y validación experimental. Las pruebas se realizaron en un entorno académico controlado y permitieron evaluar variables como la precisión de posicionamiento, la repetibilidad, el tiempo de respuesta, la capacidad de detección de objetos y el funcionamiento general del sistema. Para ello, se ejecutaron movimientos y ciclos repetitivos, se compararon las posiciones alcanzadas con las posiciones objetivo y se analizó la variación entre distintas ejecuciones.

Como resultado, se obtuvo un brazo robótico cartesiano con un movimiento más fluido, estable y repetible, capaz de ejecutar tareas automatizadas de manipulación y transferencia de productos dentro de una línea de empaque simulada. Asimismo, el sistema se desarrolló como una plataforma académica orientada a la enseñanza y experimentación en áreas relacionadas con automatización industrial, robótica, sistemas embebidos, sensores y visión computacional. Aunque el proyecto no fue validado bajo condiciones industriales de producción continua, permitió demostrar la viabilidad de las mejoras mecánicas y de la automatización implementada dentro de un entorno de laboratorio.